\documentclass[11pt,a4paper]{article}

%==================================================
% Basic setup
%==================================================
\usepackage[utf8]{inputenc}
\usepackage[T1]{fontenc}
\usepackage[english]{babel}
\usepackage{lmodern}
\usepackage{microtype}
\usepackage{geometry}
\geometry{margin=2.5cm}
\usepackage{setspace}
\onehalfspacing

%==================================================
% Mathematics, figures, tables
%==================================================
\usepackage{amsmath,amssymb,amsfonts,mathtools}
\usepackage{graphicx}
\usepackage{booktabs}
\usepackage{array}
\usepackage{enumitem}
\usepackage{caption}
\usepackage{subcaption}
\usepackage{tabularx}

%==================================================
% Links and colors
%==================================================
\usepackage{xcolor}
\usepackage{hyperref}

\hypersetup{
    colorlinks=true,
    linkcolor=blue,
    citecolor=blue,
    urlcolor=blue
}

%==================================================
% Simple pedagogical boxes
%==================================================
\newenvironment{guidancebox}
{\par\medskip\noindent\textbf{Guidance.}\begin{quote}\itshape}
{\end{quote}\medskip}

\newenvironment{claimbox}
{\par\medskip\noindent\textbf{Main claim.}\begin{quote}}
{\end{quote}\medskip}

\newenvironment{limitationbox}
{\par\medskip\noindent\textbf{Limitation.}\begin{quote}}
{\end{quote}\medskip}

%==================================================
% Project metadata
%==================================================
\title{
    \textbf{Final Project Manuscript}\\
    \textbf{Mathematical Modeling of Chocolate Toxcity for Dogs}\\[0.5em]
    \large BS-3-S2-EC-PHYMAT / Systèmes physiologiques et modélisation mathématique
}

\author{
    Team 6 -- Third-Year Biotechnology and Bioinformatics, INSA Lyon\\[0.5em]
    Olinda Mendes, Coppélia Reynaert-Verbiese, Doriane Rombaut, Capucine Schmidt, Thao Van Tran, Marion Venerucci
}

\date{Academic year 2025--2026}

\begin{document}

\maketitle

%==================================================
% Team roles
%==================================================
\begin{center}
\begin{tabular}{p{4cm}p{6cm}p{4cm}}
\toprule
\textbf{Role} & \textbf{Main responsibility} & \textbf{Student name} \\
\midrule
Physiology lead & Biological mechanism and physiological consistency & Marion Venrucci \\
Mathematics/modeling lead & Variables, equations, assumptions, analysis & Van Tran \\
Computation/simulation lead & Code, numerical simulations, figures & Capucine Schmidt \\
Critical assumptions lead & Assumptions, limitations, uncertainty & Olinda Mendes \\
Communication/report lead & Writing, structure, clarity & Doriane Rombaut \\
Coordinator & Organization, integration, deadlines & Coppélia Reynaert-Verbiese \\
\bottomrule
\end{tabular}
\end{center}

\begin{guidancebox}
Roles indicate responsibilities, not isolated tasks. Every student must
understand the whole model and be able to explain the main biological and
mathematical ideas.
\end{guidancebox}

\vspace{1em}

%==================================================
% Abstract
%==================================================
\begin{abstract}
Write 150--200 words.

Your abstract should contain:

\begin{itemize}
    \item the biological question;
    \item the modeling hypothesis;
    \item the type of mathematical model used;
    \item the main result;
    \item one important limitation.
\end{itemize}

Do not include technical details, long equations, or references in the abstract.
\end{abstract}

%\tableofcontents

%==================================================
\section{Biological and Modeling Question}
\label{sec:question}
%==================================================


\subsection{Biological context}


Chez le chien, c'est la théobromine qui est responsable de la toxicité du chocolat. En effet, 
une fois dans le système gastrique, elle est absorbée dans le sang et diffusée dans l'ogranisme,
impactant ainsi des systèmes essentiels tels que le système nerveux ou cardiovasculaire. \\
La théobromine ayant un impact sur l'entièreté de l'organisme, nous avons choisi de prendre 
comme système le chien dans son ensemble. Ce système abrite ainsi l'évolution de la concentration
de la théobromine dans le sang et son impact sur l'organisme. En effet, en augmentant, la quantité
de théobromine entraîne une réaction du corps dont découle l'apparition de certains symptômes.
C'est le lien entre cette concentration et ces signes naissants que nous cherchons à comprendre. 
Cette réaction peut être observée sur plusieurs heures. 
\\



\subsection{Modeling question}

State the precise modeling question.


\begin{quote}
\itshape
Comment est-ce que la théobromine évolue au cours du temps et agit sur l'évolution des symptômes
chez le chien ? A partir de quelle seuil la dose ingérée est-elle mortelle ? 
\end{quote}


\subsection{Reference regime and perturbation}

Pour répondre à cette question, nous prendrons comme régime de référence un modèle sain, c'est-à-dire
un chien n'ayant pas ingéré de chocolat. La perturbation étudiée est l'entrée dans le système
de chocolat, et donc de théobromine, à un instant $t{_0}$. 



%==================================================
\section{System Boundary, Variables, Parameters, and Assumptions}
\label{sec:variables-assumptions}
%==================================================

\begin{guidancebox}
This section defines the model before equations are written. A model with
unclear variables or hidden assumptions is not scientifically controlled.
\end{guidancebox}

\subsection{System boundary}

State what is inside and outside the model.

\begin{itemize}
    \item Dans le modèle :
    \begin{itemize}
        \item {Voies circulatoires et sang}
        \item {Système digestif}
        \item {Autres voies métaboliques indirectement impliquée : système nerveux, respiratoire...}
        \item {Organes}
        \item {Récepteurs, protéines, enzymes}
        \item {Hormones et neurotransmetteurs}
    \end{itemize}

    \item A l'extérieur du modèle :
    \begin{itemize}
        \item {Potentiels problèmes de santé antérieurs du chien}
        \item {Température pouvant impacter l'activité du métabolisme}
        \item {Le reste est-ce que c'est vraiment en dehors ?}
    \end{itemize}
\end{itemize}

\subsection{State variables}

\begin{center}
\begin{tabular}{p{2.5cm}p{8cm}p{3cm}}
\toprule
\textbf{Symbol} & \textbf{Meaning} & \textbf{Unit} \\
\midrule
\(Q(t)\) & {Masse de théobromine présente dans le système digestif} & {mg} \\
\(D(t)\) & {Concentration de théobromine absorbée dans le sang par la masse corporelle} & {mg/kg} \\
\(S(t)\) & {Intensité des symptômes} & {\%} \\
\bottomrule
\end{tabular}
\end{center}

\subsection{Parameters}

\begin{center}
\begin{tabularx}{\textwidth}{p{2.3cm}X p{2.5cm}}
\toprule
\textbf{Symbol} & \textbf{Meaning} & \textbf{Unit} \\
\midrule
\(M{_{choc}}\) & \emph{Masse de chocolat ingérée} & \emph{g} \\
\(C{_{type}}\) & \emph{Concentration en théobromine selon le type de chocolat} & \emph{mg/g} \\
\(W\) & \emph{Masse du chien} & \emph{kg} \\
\(f\) & \emph{Fraction de théobromine no évacuée directement par l'organisme} & \emph{\%} \\
\(k{_{abs}}\) & \emph{Taux d'absorption de théobromine dans le sang} & \emph{$h^{-1}$} \\
\(k{_{e}}\) & \emph{Taux d'élimination de la théobromine} & \emph{$h^{-1}$} \\
\(\rho\) & \emph{Taux d'activation des symptômes} & \emph{$h^{-1}$} \\
\(K{_{D}}\) & \emph{Seuil d'activation des symptômes } & \emph{mg/kg} \\
\(n\) & \emph{Coefficient de Hill} & \emph{Sans unité} \\
\(\mu{_S}\) & \emph{Taux de récupération des symptômes} & \emph{$h^{-1}$} \\

\bottomrule
\end{tabularx}
\end{center}

\subsection{Inputs and observables}
\paragraph{Inputs.}
Nos variables d'entrée sont toutes celles liées à l'ingestion de chocolat. On compte ainsi : 
\begin{itemize}
        \item {La masse de chocolat ingérée : \(M{_{choc}}\)}
        \item {La concentration en théobromine pour le type de chocolat ingéré : \(C{_{type}}\)}
        \item {La fraction de théobromine non évacuée direactement par l'organisme : \(f\)}
\end{itemize}

Describe your input protocol.

\paragraph{Observables.}
Les observables sont des points caractéristiques des courbes, apportant des précisions
sur le modèle biologique et ses limites. \\
\begin{itemize}
    \item {La quantité de théobromine présente dans le système digestif décroit linéairement :
    il n'y a pas d'observables}
    \item {La concentration en théobromine dans le sang commence d'abord par augmenter
    puis diminue. Elle atteint un maximum : \(D{_{max}}=\lim_{t \to \infty} D(t)\)}
    \item {L'intensité des symptômes évolue en suivant une courbe sigmoïdale. Elle atteint
    donc une saturation maximale : \(S{_{max}}=\lim_{t \to \infty} S(t)\)}
    \item {Certains temps peuvent être caractéristiques. On y inclut : 
    \begin{itemize}
        \item {Le temps d'atteinte du seuil d'activation des symptômes : 
        \(T{_{thr}}|S(T{_{thr}}-1)=0 \cap S(T{_{thr}})\ge 0\)}
        \item {Le temps de rétablissemnt : \(T{_{ret}}|S(T{_{thr}-1})\ge 0 \cap S(T{_{thr}})=0\)}
    \end{itemize}}
\end{itemize}

Examples:

\[
X_{\max},
\qquad
T_{\mathrm{rec}},
\qquad
\int_0^T X(t)\,dt,
\qquad
X(T).
\]

List the observables used in your project.

\subsection{Assumptions}

\begin{center}
\begin{tabular}{p{5cm}p{5cm}p{4cm}}
\toprule
\textbf{Assumption} & \textbf{Reason} & \textbf{Possible limitation} \\
\midrule
\emph{Le taux d'absorption est constant (et le même chez tous les chiens)} 
& \emph{Pour simplifier le modèle} 
& \emph{Le taux d'absorption pourrait dépendre de plusieurs facteurs 
comme la digestion/ aliments dans l'intestin (poids ? race ?)} \\
\emph{Il n'y a pas d'intervention vétérinaire pendant la simulation} 
& \emph{Pour simplifier le modèle et étudier son évolution naturelle} 
& \emph{Ne représente pas forcément une situation réelle} \\
\emph{Le taux d'élimination est constant (et le même chez tous les chiens)} 
& \emph{Pour simplifier le modèle} 
& \emph{Le taux d'élimination peut varier chez différents types de chien. (poids ?)} \\
\emph{La quantité de théobromine est la même pour tous les marques de chocolat} 
& \emph{Conserver un nombre raisonnable de comparaisons} 
& \emph{En réalité, la quantité de théobromine peut varier selon la marque du chocolat} \\
\emph{Les seuils de symptômes sont applicables pour tous les chiens} 
& \emph{Obtenir une liste exhaustive de comparaisons, tirer des conclusions groupées} 
& \emph{Les seuils de toxicité sont approximatifs et peuvent dépendre de la sensibilité individuelle.} \\
\emph{Une seule ingestion de théobromine } 
& \emph{Simplifier le modèle en supprimant une évolution de la quantité de matière ingérée} 
& \emph{L'impact d'une dose augmentée au cours du temps est ignorée de même qu'une potentielle diminution de l'efficacité du métabolisme } \\
\bottomrule
\end{tabular}
\end{center}
%==================================================
\section{Model Construction}
\label{sec:model}
%==================================================

\begin{guidancebox}
Each equation must be readable as a biological sentence. Do not write a term
unless you can explain what it represents.
\end{guidancebox}

\subsection{Conceptual diagram}

Insert your conceptual diagram.

\begin{figure}[h]
    \centering
    %\includegraphics[width=0.8\textwidth]{figures/conceptual_diagram.pdf}
    \fbox{\parbox{0.8\textwidth}{
    \centering
    Insert conceptual diagram here.
    }}
    \caption{Conceptual diagram of the modeled system. Arrows indicate
    activation, production, transfer, degradation, inhibition, or feedback.}
    \label{fig:conceptual-diagram}
\end{figure}

\subsection{Model equations}

Write the final model.

For example:

\[
\frac{dD}{dt}
=
\text{absorbtion}
-
\text{élimination}.
\]

Then give the actual equations:

\[
\frac{dQ}{dt}=-k{_{abs}}Q,
\]

\[
\frac{dD}{dt}=\frac{k{_{abs}}}{W}Q-k{_{e}}D,
\]

\[
\frac{dS}{dt}=\rho\frac{D^n}{K{_D}^n + D^n}(1-S)-\mu{_S}S.
\]


\subsection{Meaning of each term}

Explain each term.

For example:

\[
-k{_{abs}}Q
\]

réprésente la sortie de théobromine du système digestif vers le sang avec un taux d'absorption
\(k{_{abs}} \ge 0\).

\[
\frac{k{_{abs}}}{W}Q
\]

répresente la concentration de théobromine absorbée dans le sang depuis le système digestif, 
en fonction du poids du chien. 

\[
-k{_{e}}D
\]

représente l'élimination de la théobromine du sang avec un taux d'élimination \(k{_{e}} \ge 0\).

\[
\rho\frac{D^n}{K{_D}^n + D^n}
\]

représente l'évolution des symptômes avec saturation. Une fonction de Hill est utilisée de 
manière à encadrer la variable avec deux limites. La constante $\rho$ fait évoluer ce terme 
à un certain taux. 

\[
(1-S)
\]

permet de maintenir S entre 0 et 1. 

\[
-\mu{_S}S
\]

représente la diminution des symptômes, c'est-à-dire le rétablissement du chien, à un taux
$\mu{_S}$

\subsection{Units and dimensional consistency}

\[
\frac{dQ}{dt}=\frac{masse (mg)}{temps (h)},
\]

et

\[
-k{_{abs}}Q=temps^{-1} (h^{-1}) \times masse (mg),
\]

donc les dimensions sont cohérentes. 

\[
\frac{dD}{dt}=\frac{concentration (mg/kg)}{temps (h)}
\]

et 

\[
\frac{k{_{abs}}}{W}Q-k{_{e}}D=\frac{temps^{-1} (h^{-1})}{masse (kg)} \times masse (mg)-temps^{-1} (h^{-1})
\times concentration (mg/kg),
\]

Les dimensions sont cohérentes.

\[
\frac{dS}{dt}=\frac{intensité (sans unité)}{temps (h)}.
\]

et

\[
\rho\frac{D^n}{K{_D}^n + D^n}(1-S)-\mu{_S}S=\rho\frac{D^n}{K{_D}^n + D^n}(1-S)-\mu{_S}S.
\]

\subsection{Biological domain}

State the domain of the variables.

Examples:

\[
X(t)\geq 0,
\]

\[
0\leq B(t)\leq 1,
\]

\[
0\leq V(t)\leq 1.
\]

Explain whether the model preserves this domain.

%==================================================
\section{Mathematical Analysis}
\label{sec:analysis}
%==================================================

\begin{guidancebox}
This is the mathematical core of the report. It should contain more than
simulation. At minimum, identify equilibria and discuss stability.
\end{guidancebox}

\subsection{Equilibria}

Set the derivatives equal to zero.

For a one-variable model:

\[
\frac{dX}{dt}=f(X),
\]

equilibria satisfy

\[
f(X^*)=0.
\]

For a two-variable model:

\[
\frac{dX}{dt}=f(X,Y),
\qquad
\frac{dY}{dt}=g(X,Y),
\]

equilibria satisfy

\[
f(X^*,Y^*)=0,
\qquad
g(X^*,Y^*)=0.
\]

List the equilibria and interpret them biologically.

\subsection{Stability}

For a one-dimensional model, compute

\[
f'(X^*).
\]

If

\[
f'(X^*)<0,
\]

the equilibrium is locally stable.

If

\[
f'(X^*)>0,
\]

the equilibrium is unstable.

For a two-dimensional model, compute the Jacobian:

\[
J(X,Y)
=
\begin{pmatrix}
\frac{\partial f}{\partial X} & \frac{\partial f}{\partial Y}
\\[0.4em]
\frac{\partial g}{\partial X} & \frac{\partial g}{\partial Y}
\end{pmatrix}.
\]

Evaluate \(J\) at equilibrium and discuss trace, determinant, or eigenvalues.

\subsection{Time scales}

Identify the main time scales.

Examples:

\[
\tau=\frac{1}{k},
\]

or

\[
t_{1/2}=\frac{\ln 2}{k}.
\]

Interpret them biologically.

\subsection{Thresholds, feedback, or bifurcations}

Use this subsection only if relevant.

State clearly whether the project involves:

\begin{itemize}
    \item an operational threshold;
    \item a dynamical threshold;
    \item positive feedback;
    \item negative feedback;
    \item bistability;
    \item hysteresis;
    \item bifurcation;
    \item irreversible loss.
\end{itemize}

\begin{limitationbox}
Do not claim bistability, hysteresis, Hopf bifurcation, or point of no return
unless your equations actually support that interpretation.
\end{limitationbox}

%==================================================
\section{Simulation Protocol}
\label{sec:simulation-protocol}
%==================================================

\begin{guidancebox}
A simulation is an experiment on the model. It must be reproducible.
\end{guidancebox}

\subsection{Parameter values}

Give the parameter set used for baseline simulations.

\begin{center}
\begin{tabular}{p{2.5cm}p{3cm}p{3cm}p{5cm}}
\toprule
\textbf{Parameter} & \textbf{Value} & \textbf{Unit} & \textbf{Justification} \\
\midrule
\(k\) &  &  &  \\
\(\alpha\) &  &  &  \\
\(K\) &  &  &  \\
\(n\) &  &  &  \\
\bottomrule
\end{tabular}
\end{center}

\subsection{Initial conditions}

State the initial conditions:

\[
X(0)=X_0,
\qquad
Y(0)=Y_0,
\qquad
Z(0)=Z_0.
\]

Explain their biological meaning.

\subsection{Input protocol}

Describe any external input.

Example:

\[
u(t)=
\begin{cases}
u_0, & 0\leq t\leq T_{\mathrm{stim}},\\
0, & t>T_{\mathrm{stim}}.
\end{cases}
\]

\subsection{Numerical method}

State the software and solver used.

Example:

\begin{quote}
The ODE system was integrated in Python using \texttt{solve\_ivp} from
\texttt{scipy.integrate}, with relative tolerance \(\ldots\) and absolute
tolerance \(\ldots\).
\end{quote}

or

\begin{quote}
The ODE system was integrated in MATLAB using \texttt{ode45}.
\end{quote}

\subsection{Simulation observables}

Define the quantities extracted from simulations.

Examples:

\[
X_{\max}
=
\max_{0\leq t\leq T}X(t),
\]

\[
\mathrm{AUC}
=
\int_0^T X(t)\,dt,
\]

\[
T_{\mathrm{rec}}
=
\inf
\left\{
t\geq 0:
|X(s)-X^*|<\varepsilon
\text{ for all }s\geq t
\right\}.
\]

%==================================================
\section{Results}
\label{sec:results}
%==================================================

\begin{guidancebox}
Each figure must answer a question. Do not include plots only because they look
nice.
\end{guidancebox}

\subsection{Baseline dynamics}

Present the baseline simulation.

\begin{figure}[h]
    \centering
    %\includegraphics[width=0.75\textwidth]{figures/baseline_dynamics.pdf}
    \fbox{\parbox{0.75\textwidth}{
    \centering
    Insert baseline trajectory figure here.
    }}
    \caption{Baseline dynamics of the model. The figure should show the main
    state variables and the relevant time scale.}
    \label{fig:baseline}
\end{figure}

Interpret the figure.

Questions:

\begin{itemize}
    \item Does the system approach equilibrium?
    \item Is the response transient or persistent?
    \item Is there overshoot?
    \item What is the recovery time?
    \item Does the result match the stability analysis?
\end{itemize}

\subsection{Perturbation or treatment scenario}

Present the main perturbation scenario.

\begin{figure}[h]
    \centering
    %\includegraphics[width=0.75\textwidth]{figures/perturbation_scenario.pdf}
    \fbox{\parbox{0.75\textwidth}{
    \centering
    Insert perturbation or treatment figure here.
    }}
    \caption{Effect of the selected perturbation on model dynamics.}
    \label{fig:perturbation}
\end{figure}

Interpret the result.

\subsection{Scenario comparison}

Compare at least two scenarios.

Examples:

\begin{itemize}
    \item low versus high exposure;
    \item weak versus strong feedback;
    \item fast versus slow clearance;
    \item early versus late treatment;
    \item different initial conditions.
\end{itemize}

%==================================================
\section{Sensitivity and Robustness Analysis}
\label{sec:sensitivity}
%==================================================

\begin{guidancebox}
This section shows whether your conclusion depends on one arbitrary parameter
choice or remains true across a reasonable range.
\end{guidancebox}

\subsection{One-parameter sweep}

Choose one important parameter \(\lambda\) and vary it.

Define an observable \(Q(\lambda)\).

Examples:

\[
k \longmapsto T_{\mathrm{rec}}(k),
\]

\[
D \longmapsto X_{\max}(D),
\]

\[
\alpha \longmapsto X(T;\alpha).
\]

\begin{figure}[h]
    \centering
    %\includegraphics[width=0.75\textwidth]{figures/parameter_sweep.pdf}
    \fbox{\parbox{0.75\textwidth}{
    \centering
    Insert parameter sweep figure here.
    }}
    \caption{Effect of parameter variation on the selected observable.}
    \label{fig:parameter-sweep}
\end{figure}

\subsection{Robustness of the main conclusion}

State whether the main conclusion is robust.

Examples:

\begin{quote}
The conclusion remains valid when \(k\) varies over the interval
\([k_{\min},k_{\max}]\).
\end{quote}

or

\begin{quote}
The conclusion depends strongly on the feedback strength \(\alpha\), so the
result should be interpreted cautiously.
\end{quote}

%==================================================
\section{Discussion and Model Criticism}
\label{sec:discussion}
%==================================================

\begin{guidancebox}
This section is not a place to repeat the results. It is where you explain what
the model suggests and what it cannot prove.
\end{guidancebox}

\subsection{Main biological interpretation}

Explain the biological meaning of the results.

\begin{claimbox}
State your main claim in one or two sentences.

Example: The model suggests that recovery time is controlled primarily by the
clearance rate, while peak response is controlled primarily by the input
amplitude.
\end{claimbox}

\subsection{What the model supports}

List the conclusions that are directly supported by your analysis and
simulations.

\begin{itemize}
    \item {[Supported conclusion 1]}
    \item {[Supported conclusion 2]}
    \item {[Supported conclusion 3]}
\end{itemize}

\subsection{What the model does not prove}

List conclusions that would be overclaims.

\begin{itemize}
    \item {[Overclaim 1]}
    \item {[Overclaim 2]}
\end{itemize}

\subsection{Limitations}

Discuss limitations.

Possible limitations:

\begin{itemize}
    \item simplified physiology;
    \item missing spatial structure;
    \item no delay;
    \item constant parameters;
    \item uncertain parameter values;
    \item no individual variability;
    \item no direct data fitting;
    \item phenomenological variables.
\end{itemize}

\subsection{Possible improvements}

State how the model could be improved.

Examples:

\begin{itemize}
    \item adding a second compartment;
    \item adding a delay;
    \item using experimental data to estimate parameters;
    \item adding individual variability;
    \item testing an alternative feedback structure.
\end{itemize}

%==================================================
\section{Conclusion}
\label{sec:conclusion}
%==================================================

Write a short conclusion.

It should answer:

\begin{itemize}
    \item What was the biological question?
    \item What was the model?
    \item What was the main mathematical result?
    \item What is the main biological interpretation?
    \item What is the main limitation?
\end{itemize}

A good conclusion is precise and modest.

Do not introduce new results here.

%==================================================
\section{AI and Software Use Statement}
\label{sec:ai-software}
%==================================================

\subsection{AI use}

State how AI tools were used.

Examples:

\begin{itemize}
    \item brainstorming biological mechanisms;
    \item identifying hidden assumptions;
    \item checking clarity of explanations;
    \item debugging code;
    \item improving English formulation.
\end{itemize}

State clearly what was verified by the group.

\begin{center}
\begin{tabular}{p{2cm}p{4cm}p{5cm}p{3cm}}
\toprule
\textbf{Date} & \textbf{Task} & \textbf{Prompt or use} & \textbf{Group decision} \\
\midrule
 & & & \\
 & & & \\
 & & & \\
\bottomrule
\end{tabular}
\end{center}

\subsection{Software use}

State the software used.

Examples:

\begin{itemize}
    \item Python;
    \item MATLAB;
    \item Jupyter Notebook;
    \item NumPy/SciPy;
    \item Matplotlib;
    \item Git or shared folder.
\end{itemize}

%==================================================
\section*{Acknowledgments}
\addcontentsline{toc}{section}{Acknowledgments}
%==================================================

Optional.

Mention help received from instructors, classmates, or external sources if
relevant.

%==================================================
\begin{thebibliography}{99}
\addcontentsline{toc}{section}{References}

\bibitem{Strogatz}
S. H. Strogatz,
\textit{Nonlinear Dynamics and Chaos: With Applications to Physics, Biology,
Chemistry, and Engineering}.
CRC Press.

\bibitem{EdelsteinKeshet}
L. Edelstein-Keshet,
\textit{Mathematical Models in Biology}.
SIAM Classics in Applied Mathematics.

\bibitem{Murray}
J. D. Murray,
\textit{Mathematical Biology I: An Introduction}.
Springer.

\bibitem{Ingalls}
B. P. Ingalls,
\textit{Mathematical Modeling in Systems Biology: An Introduction}.
MIT Press.

\bibitem{Alon}
U. Alon,
\textit{An Introduction to Systems Biology: Design Principles of Biological
Circuits}.
Chapman and Hall/CRC.

\bibitem{Saltelli}
A. Saltelli et al.,
\textit{Global Sensitivity Analysis: The Primer}.
Wiley.

% Add project-specific references here.

\end{thebibliography}

%==================================================
\appendix

\section{Appendix A: Additional Calculations}

Include long algebraic derivations, equilibrium calculations, nondimensionalized
forms, or stability calculations that would interrupt the main text.

\section{Appendix B: Code Availability}

State where the code is available.

Example:

\begin{quote}
The Python scripts and notebooks used for the simulations are included in the
Moodle submission as supplementary files.
\end{quote}

\section{Appendix C: Additional Figures}

Include additional simulations or parameter sweeps if needed.

\end{document}